\section{Condensados de las cuatro rutas de tratamiento (INT-37)}
\label{sec:sd4-drenaje-int37}
Se selecciona drenaje gravitacional individual para cada serpentín de aire limpio situado antes de recalentamiento y ventiladores: cuatro bandejas primarias, cuatro sifones registrables y cuatro descargas visibles con ruptura de conexión. Las rutas A/B y sus reservas no comparten sifón ni una tubería presurizada de aire. La descarga se realiza hacia receptores exteriores accesibles, sin retorno a banco ni conexión con extracción H$_2$. El criterio de un sifón por sección y distinción entre presión positiva/negativa se apoya en el manual Trane, utilizado como referencia de ingeniería del drenaje; no supone seleccionar una UTA Trane ni aplicar su certificación a una fabricación propia \parencite[pp.~54--55, figs.~50--51]{int37-trane-drenaje}.

El contraste psicrométrico vigente de 900 m$^3$/h a salida ideal saturada 7,3925 grados C produce 22,203622 kg/h por ruta. Con densidad propia conservadora de 995 kg/m$^3$, son 22,315198 L/h. Se reserva factor propio cinco para descarga/limpieza/transitorio: 111,575990 L/h por ruta y 446,303960 L/h para solape de las cuatro, sin mezclar sus cierres hidráulicos. El factor no acredita fuga ilimitada de agua del serpentín ni capacidad del alcantarillado receptor. Drenaje de condensado, vaciado de circuito de agua y fuga de tubería tienen destinos y dimensionamientos separados; una fuga obliga aislamiento de rama y contención independiente.

Se fija diámetro interior mínimo de obra de 25 mm y pendiente continua no menor a 1\,\% en cada ramal aguas abajo de su sifón, sin contrapendientes ni descarga sumergida. Como cribado hidráulico a medio llenado se aplica $A=\pi D^2/8$, $R_h=D/4$, $Q=(1/n)AR_h^{2/3}S^{1/2}$, con $n=0,013$ propio. Da 230,614032 L/h, 2,067 veces la reserva por ruta. La ecuación de Manning proporciona un contraste de flujo uniforme; no demuestra paso efectivo de gotas, régimen laminar, pérdidas de sifón, biofilm o taponamiento. El coeficiente es hipótesis de obra y no ficha del material. Se seleccionan material/conexiones no combustibles resistentes al agua y al ambiente marino antes de fabricar, cotejando diámetro interior real y puerto del serpentín; no se reduce un puerto mayor para ajustarlo al cálculo \parencite[ec.~3-23 y sec.~5.5.2]{int37-epa-manning}.

Se elige cámara de serpentín en aspiración de los ventiladores, manteniendo su régimen negativo durante operación. Se fija límite propio provisional de presión negativa de 3.000 Pa en bandeja, aún por contrastar con curva de bloqueo, máxima consigna, filtro sucio y estados de reserva. Para $\rho=995$ kg/m$^3$ y $g=9,80665$ m/s$^2$, la columna equivalente es 307,452125 mm. La figura 50 de referencia exige $H=h_p+25,4$ mm, $J=H/2$ y $L=H+J+D_{ext}$; se eligen $H=350$ mm, $J=175$ mm y diámetro exterior de cálculo 32 mm, dando $L=557$ mm. Se reserva espacio vertical libre mínimo de 650 mm bajo el puerto, antes de tolerancias/adaptadores. La altura H reserva 17,147875 mm sobre la fórmula mínima, y la envolvente vertical deja 93 mm para herrajes/tolerancias bajo diámetro exterior supuesto. Un diámetro exterior mayor o presión que exceda la cota exige redimensionar y ampliar soporte; no basta el caudal nominal del ventilador para justificar esta cota. H/J/L identifican las alturas de esa figura y no dimensiones intercambiables de un sello positivo \parencite[p.~55, fig.~50]{int37-trane-drenaje}.

Cada sifón tiene acceso propio para cebado y limpieza; se llena antes de energizar y después de una parada que pueda secarlo. La descarga termina con separación abierta propia mínima 50 mm sobre receptor no inundable y con evacuación gravitacional accesible; es reserva de proyecto, su cumplimiento sanitario se coteja en ingeniería de instalación. No se conecta venteo abierto entre bandeja y sello que eluda la columna hidráulica. Las inspecciones comprueban sello, flujo y contrapresión; un sifón seco o pérdida de presión invalida la disponibilidad de la ruta. El sello de agua no es válvula certificada contra H$_2$, y no resuelve retorno de gas cuando fallan compuertas/ventilación. La presión positiva requiere geometría distinta: el presente esquema negativo no se presenta como apto para inversión de presión.

Cada ruta incorpora contención secundaria independiente y nivel alto con alarma/retirada de agua a través de su cadena local. Se reserva bandeja secundaria de obra de 1,0 por 0,8 por 0,10 m, con nivel útil máximo 0,05 m, 40 L antes de la reserva superior. Sin intrusiones ni aporte de fugas, el cribado da 21,510004 minutos al caudal propio de 111,575990 L/h; no es tiempo autorizado de respuesta ni contención de fuga de circuito. El serpentín y la bandeja primaria reales pueden exigir ampliar la superficie; la bandeja secundaria no reemplaza su selección ni oculta sobrepaso. Interruptor de nivel, interfaz de retirada y detección de señal inválida se seleccionan y ensayan antes de habilitación: su mera especificación no completa el enclavamiento.

La reserva vertical se integra con soportes, plano, pesos, aislamiento INT-34 y mantenimiento sin ocupar pasillos ni sumar una conexión gaseosa entre rutas. La recepción prevista prueba descarga con ventiladores a máxima presión adoptada, filtro sucio, transferencia de rutas, cebado, obstrucción, rebose y retirada; registra caudal, presión real de bandeja, nivel y salida segura. El diseño selecciona drenaje y desarrolla cantidades/dimensiones de obra. Conserva pendientes la presión máxima real, diámetro/puerto/material seleccionados, receptor sanitario, conjunto húmedo, soporte/implantación y alarma independiente.

Los sifones y descargas abiertas no son barreras de humo/fuego ni impiden difusión de gas cuando se secan o se invierte la presión. La transferencia y la secuencia gas/fuego deben mantener la cámara limpia separada del banco y comprobar presión/sello en los estados de fallo; una inversión invalida esta ruta hasta resolver su aislamiento. No se autoriza sellado/extinción automática de bancos mediante este drenaje ni se consideran aprobadas sus penetraciones por disponer sello de agua.
